\noindent{\large\bfseries 10 月 21 日补充：执行规则、未成交原因与原账本接入}
\par\noindent{\small 前文保留 10 月 19 日的原实验与正文；本补充单独说明新增代码和实际运行范围。}

\lead{先说明这次究竟接到了什么}
新提供的 \texttt{H20\_equal\_G1.zip} 实际覆盖 sample1：2007 年 2 月 2 日至 2009 年 12 月 31 日，共 711 个交易日；不是建议先接的 sample3 或 2024 年窗口。10 月 21 日 15:08 通过原有授权只读接口查询一行，确认数据库目前最早的 L2 为 2024 年 1 月 2 日 09:15。因此两者没有日期交集，不能把旧因子交易搬到 2026 年行情上回放。

包内 \texttt{orders.csv} 有 681 行，但每行是某日执行尝试的汇总，没有逐证券订单标识、方向和股数。成交载荷有 30,972 行，保存买卖金额、费用和成交后复权持仓单位；其中 \texttt{units} 不是成交股数。账本清单所列原始开盘价、复权开盘价、复权因子实际放在包外共享目录，没有随本包给出。持仓事件还需要缺失的证券字典才能恢复证券身份和估值。文件名中有 ``orders'' 或 ``trades''，不等于已经具备真实券单与原始股数。

由首日现金与当天买卖、费用推得期初归一化现金约为 1，逐日累加卖出金额、减买入金额和费用，711 日现金与原表的最大差约为 $4.55\times10^{-15}$。这是现金核对；缺少独立期初账户声明，也没有独立还原逐股持仓市值，更没有把归一化单位直接解释为人民币或股票股数。原表有 703 日报告持有未取得当日报价的证券，最多 32 只，其估值需额外核实。将固定实际成交路径中的费用退回现金，也只是一项核算，不等于重新生成零费策略或完成了流动性回放。

\begin{table}[H]\centering\small
\begin{tabularx}{\linewidth}{p{30mm}Y}
\toprule 最小补充资料 & 为什么需要\\\midrule
有交集的原组合 & 同组合 sample3 中一段 2024 年的原始资料；按时间和数据完整性选窗，不按收益选窗。\\
逐证券券单 & 订单关联、证券、信号及执行时间、方向、股数或目标权重、参考价和覆盖/撤单语义。\\
完整原成交与共享数据 & 原始成交价、数量、费用及订单关联；账本引用的价格、复权因子、证券字典和实际读取实现。只有日级时间时保留日级精度。\\
窗口前账户与基准 & 现金、真实持仓、可卖量、负债、应收和锁定股份；同期每日净值/持仓，以及 100001 价格指数。\\
\bottomrule
\end{tabularx}
\caption{只需一个组合的小窗口。私人因子明细和完整账户继续留在授权本地。}
\end{table}

\lead{旧结果的 22 笔不能混为一种未完成}
1 亿元、10\% 的旧场景共 93 笔需求，其中 22 笔当日未精确完成：14 笔较大余量主要由参与率额度不足解释；7 笔是不足 100 股尾差；1 笔是 000635.SZ 在 8 月 31 日的明确停牌卖单。7 笔尾差合计 504 股、参考金额 10,463.86 元，其中有历史尾差在之后调仓再次出现。停牌单的 10,300 股在 9 月 1 日 09:30:06 完成。

\begin{table}[H]\centering\small
\begin{tabularx}{\linewidth}{p{28mm}Y}
\toprule 分类 & 证据与解释\\\midrule
正常市场障碍 & 明确停牌；或官方方向性限价与原始盘口无对手盘同时成立。模型扣减后的零深度不能算原始市场无对手盘。\\
已设限额不足 & 市场正常交易，但已定参与率、执行时段及有效来源量给出的额度小于需求；不能自动称为程序故障。\\
程序或模型限制 & 逐快照取整损失、有效子单余量取消、同价深度全天累计扣减等；保留具体规则与影响。\\
数据问题或待查 & 缺行情、区间量额不一致、身份和状态证据不足；不能自动填成正常或停牌。\\
\bottomrule
\end{tabularx}
\caption{正常、异常和参数受限分别解释。小尾差保留精确数量，未成交仍留在全部需求分母。}
\end{table}

\clearpage
\lead{新规则怎样使用已经发生的成交量}
旧规则对每张快照计算 $100\lfloor pV/100\rfloor$，零碎额度不能累计。新增累计模式只积累存在活动需求时、已经观察到且质量检查有效的已结束区间量：先以 $pV$ 记入已知额度，再按申报规则从可用余额中发子单。不是预先使用全天成交量；当前尚未结束区间和未来行情不增加先前额度。

同一股票的全部买卖单共享额度，按绝对数量占用，不能把买卖相抵。待执行子单先预占，成交转为已用，取消或到期释放未成交预占；午休、跨日、未知间隔和来源异常处清空可用余额并留下记录，不把旧时段额度搬到不连续时段。金额与成交量原字段保留，负数和缺失不被无痕填零。

这一模式改变了旧版“仅前一区间额度”的规则。旧默认与旧结果保留用于对照，新增实验显式记录模式选择。预占、成交、释放和丢弃均写入额度流水，以账户等式和逐事件数量检查是否重复使用。

\lead{盘口与部分成交采用什么假设}
原保守模式将同一股票、方向、价格当天的已模拟成交量持续扣除。这一压力假设保留，但原始盘口数量、模拟已消耗数量和模型剩余深度分别记录。旧 22 单诊断中已有 219 次原始十档有量、模型扣后为零的情形，原因修正为模型深度耗尽，不能再称真实市场无对手盘。

新版比较情景使用每张盘口刷新可见数量，并让同股同一快照的多单共享该盘口。快照不能重建完整委托队列，因此刷新只是明确的比较假设；更高完成率不代表实盘保证成交，也没有据此识别永久市场冲击。

已合法发出的子单若部分成交，剩余部分在相同有效时段和数据连续条件下继续保留同一个子单身份；新发买单仍须符合整手规则。卖出余股与买入新单分开处理。新增规则仅适用这次 \texttt{000xxx.SZ} 样本，没有把它套用到全部交易所与板块。

深交所规则要求买入按 100 股或整数倍申报；非整百持股的全部余数可一次单独卖出，或与整百部分合并申报，不能任意拆零。有效子单部分成交后留下 52、99 或 51 股，与新发 52、99 或 51 股买单是不同情况。依据为\href{https://www.szse.cn/lawrules/rule/trade/current/t20150914_565054.html}{深交所卖出申报数量通知}及\href{https://docs.static.szse.cn/www/lawrules/rule/trade/current/W020260424690713155663.pdf}{2026 年交易规则第 3.3.8 条}。

\clearpage

\lead{固定同一批母单后，究竟减少了什么}

本次固定旧 22 笔母单、原日初持仓、行情和当日终点，六种规则共实际运行 132 次；另外将同日股票同时推进，44 个订单检查与单独运行的数量、金额一致。旧 22 笔的基线数量和金额全部精确重现。诊断统一使用 1 亿元起始现金和原持仓，使资金充足，不把这一现金假设冒充原账户边界。每次另由成交、借入归还和原始末价重建现金、负债、可卖量、持仓与 NAV，最大独立 NAV 差小于 0.000001 元。

\begin{table}[H]\centering\small

\begin{tabular}{lrrr}

\toprule 规则 &14 大单余量（万股）&7 尾差余量（股）&22 单余额（万元）\\\midrule

旧规则 &710.06&504&3,864.23\\

仅累计额度 &635.04&403&3,307.21\\

仅刷新盘口 &704.46&353&3,784.10\\

仅保留子单余量 &706.40&401&3,811.24\\

仅改卖出余股规则 &710.06&207&3,863.93\\

新比较情景 &630.50&104&3,250.27\\

\bottomrule\end{tabular}

\caption{全部仍使用 10\% 参与率。余额是历史需求剩余的参考金额合计，不是收益损失或期末挂单。}\end{table}

新比较情景下，14 笔较大订单仍全部未精确完成，剩余 6,305,000 股、参考金额 32,432,858 元。7 笔尾差中 5 笔精确完成，余下 104 股、2,314 元，是 000048 在 8 月 24、31 日的两笔原母单各自带入 52 股；固定券单不能凭修改数量规则虚构这两笔买入。明确停牌的 10,300 股不变。

\begin{table}[H]\centering\small\begin{tabular}{lrrr}

\toprule 新比较情景的结果 &笔数&余量（股）&余额（元）\\\midrule

已设限额不足 &14&6,305,000&32,432,858.00\\

精确完成 &5&0&0.00\\

数量规则与历史尾差 &2&104&2,314.00\\

正常市场障碍：明确停牌 &1&10,300&67,568.00\\

\bottomrule\end{tabular}

\caption{数据质量与数量原因并列保留，不能因归入额度不足就抹去区间缺量。}\end{table}

\lead{三个可手算例子：绝对几百万元，不等于额度足够}

\begin{table}[H]\centering\small\begin{tabular}{lrrrrr}

\toprule 股票与日期 &需求&全日市场量&旧取整额度&旧成交&新成交\\\midrule

000529 08-03 &141.84&571.56&49.26&49.26&54.06\\

000550 08-03 &39.19&325.59&22.70&20.34&31.67\\

000530 08-24 &127.42&1,305.85&120.31&120.31&125.86\\

\bottomrule\end{tabular}

\caption{数量单位均为万股。新规则还受边界、预占和盘口约束，不等于有效窗口总额度全能成交。}\end{table}

000529 的需求为 1,418,400 股，而全日成交量的 10\% 仅 571,560 股；有效窗口未取整额度更少，为 545,211.6 股。000550 的有效额度为 318,200.1 股，旧逐张整手取整仅 227,000 股，旧模型实际成交 203,400 股；新情景成交 316,700 股，仍少于 391,900 股需求。000530 的全日 10\% 为 1,305,849.8 股，看似超过需求，但有效窗口未取整仅 1,261,339.8 股；旧取整后 1,203,100 股，新成交 1,258,600 股，仍剩 15,600 股。

上述有效未取整额度沿用旧口径：该来源之后还须有同一时段、无未知间隔的可执行快照。新增累计模式的已观察额度还可能包含随后遇到边界而必须丢弃的部分，不能全当可成交。三例累计记入分别为 552,541.6、318,760.1、1,262,949.8 股，实际成交后分别丢弃 11,941.6、2,060.1、4,349.8 股；各额度段最终预占均为零。两种来源口径在 CSV 中分列，不混作同一个上限。

复用此前 150 个股票日的量额检查：149 组日末累计与日线一致、1 组明确停牌；45 组区间求和不足，合计少 404,200 股、9,342,869 元。没有重扫同一套检查，也没有用事后累计差分覆盖原字段。L2 价格除以 10,000 为元/股，成交量与盘口量为股，成交额直接为元；逐笔金额按数量乘价格核对。日线与 L2 同源，日末对齐不构成盘中完整性的独立证明。

\clearpage

\lead{连续账户重新运行后的收益与完成率}

从 7 月 31 日同样的 1 亿元现金开始，新增三条完整 26 日路径：参考假设、L2 不限参与率、L2 的 10\% 参与率。现金、持仓、可卖量和债务逐日保留；以后需求按各账户实际状态重算。旧 25 路径结果保留，以下旧 10\% 行为原已运行结果，新设置均重新执行全期，未在旧收益上手工加减。

\begin{table}[H]\centering\small\begin{tabular}{lrrr}

\toprule 连续账户 &全期收益（\%）&当日完成（\%）&对参考收益差（百分点）\\\midrule

按参考价全成假设 &-1.3017&100.0000&+0.0000\\

旧规则 10\% &-3.2361&87.6096&-1.9343\\

新比较情景 10\% &-2.9764&89.6872&-1.6746\\

新比较情景不限参与率 &-4.2798&99.9759&-2.9781\\

\bottomrule\end{tabular}

\caption{本表是调试动量的连续策略对照；参考价全成不是原回测实际执行，也没有同期 100001 指数超额。}\end{table}

新 10\% 情景比旧规则的全期收益变化 +0.2597 个百分点，期末净资产变化 +259,707.17 元；相对参考假设的净资产差为 -1,674,634.50 元。这里的差额属于规则与后续持仓路径的比较，不能冒称正式因子相对原回测执行的流动性差。完成率较高也不保证收益较高。

新 10\% 路径共 91 笔需求，当日精确完成 76 笔；14 笔普通余量由已设限额不足解释，余 6,227,600 股、参考金额 32,054,032 元；另 1 笔明确停牌，9,900 股、64,944 元，次日续成。没有不足 100 股尾差。其全部需求金额为 311,446,782 元，未成需求保留在分母。新路径的停牌数量与旧固定单的 10,300 股不同，正是连续账户重算后续需求的结果。期末现金为零、借款 1,008,314.50 元计负债，应收和费用均为零；无仍有效挂单，历史已到期余量并未消失为假成交。

\panel{05_execution_update}{左图固定同一批 14 笔较大需求，右图为各账户重算后续需求的连续净值；两类对照的含义不同。}

代码全套 180 项测试通过。固定对照与新连续路径均通过账户核对，并另用成交、借贷和日线收盘价重建每日现金、负债、持仓与 NAV。真实样本中两笔公司行动仍为登记时零持股，非零分红和送转依靠合成测试；没有把零持股样本扩大解释为已完成真实非零权益验证。1 毫秒延迟仅为正延迟的工程设置，不代表原始行情具有 1 毫秒精度。


\clearpage
\lead{正式原券单回放与连续策略实验分开}
固定原券单模式保持给定股数和关联，不按回放账户的新净资产偷偷换单。原计划后续卖出量若超过实际买到的可卖数量，记录缺口，不超卖、不补假库存。输入若为目标权重，必须先明确缩放资金、参考价及净需求规则，不能声称各情景的股数完全相同。

前文及本次连续策略对照仍是同一动量信号、各账户按自身净资产与实际持仓重新生成需求；只是程序与模型验证。它不能替代低因子 G1、H=20 交易日、历史 100001 股票池的正式因子策略。

同一期初净资产 $A_0$、同一终点估值、两侧零费用时，正式指标定义为：
\begin{align*}
R_o&=A_o/A_0-1,&R_l&=A_l/A_0-1,&R_i&=I_T/I_0-1,\\
E_o&=R_o-R_i,&E_l&=R_l-R_i,\\
\Delta R_{\rm pp}&=100(R_l-R_o),&\Delta A&=A_l-A_o\quad\text{（元）}.
\end{align*}
其中 $o$ 为原回测实际成交的零费重建账户，$l$ 为流动性回放账户，$i$ 为同一个指数。相对净值的复利收益若另报，需另命名；多年 CAGR 不能与几周累计收益相减。原回测买 0.015\%、卖 0.065\% 的旧费用不能混入当前差额，也不能只给旧 CAGR 加回费率。

\lead{真实已运行范围与交付}
本补充包含实际代码修改、原包字段与现金核对、旧固定券单的单变量诊断及连续账户重跑。真实因子流动性回放、相对 100001 指数的超额和相对原执行的收益差，因原券单、原始价格/股数、完整账户和同期行情仍不足而未完成，正式指标不填造数。

费用、税费、最低费用、融资利息均为零，不预留现金；借款本金始终计负债。现金、股数、可卖量和权益连续承接，普通订单当日结束、明确市场障碍例外续单；未知缺失停止。30/60 分钟只是观察点，没有新增完成率或资金容量合格线。未成参考金额的历史合计既不是期末挂单，也不是收益损失。

代码和必要派生结果分别保存在内网代码仓库与本公开报告包。原始 L2、私人因子、完整持仓账户和服务器配置不进入公开包。本次没有新增或运行哈希、校验和或指纹去重流程；必要核对使用行数、具体字段、账户等式和手算案例。图源 CSV、重画脚本、复现命令及输入限制随包提供。
